\cotefile{157-4}
\scanpage{288}
\mspage{16}
% archive sheet number: 287

des types $\mathrm{Hot}_{\unsure{W|w}}$. Si on tient compte de la description de $I_W$ comme $I\times e_W$, la formule d'adjonction \unsure{inhabituelle} (2.8.3) équivaut, bien sûr, à la formule (2.8.4) -- elle s'en déduit en ``\unsure{élevant} à la puissance $I$''.

On peut renverser la vapeur, et dans une catégorie \unsure{générale} $\mathcal H$, \uline{définir} un foncteur $\pi_0:\mathcal H\longrightarrow\mathrm{Ens}$ par la formule (2.8.\deleted{\unsure{4}}5)
\[
\pi_0(X)\overset{\mathrm{def}}{=}\mathrm{Hom}(e_{\mathcal H},X)
\quad\inserted{(=\Gamma_{\mathcal H}(X),\ \text{en une notation familière})}
\]
où $e_{\mathcal H}$ est \deleted{\unsure{l'}}\inserted{un} objet final de $\mathcal H$. On aura donc, en posant $I_{\mathcal H}=I\times e_{\mathcal H}$, la formule
\[
\mathrm{Hom}(I_{\mathcal H},X)\overset{\sim}{\longrightarrow}\mathrm{Hom}(\deleted{\unsure{\pi_0}}I,\deleted{\unsure{I}}\pi_0(X)),
\]
(i.e. la formule (2.8.3)), disant que l'on a un couple $(I_{\mathcal H},\pi_0)$ de foncteurs adjoints.

Supposons d'autre part que pour tout ens. $I$, l'application canonique (d'adjonction)
\[
 I\longrightarrow\deleted{\unsure{\mathrm{Hom}}}\pi_0(I_{\mathcal H})=\mathrm{Hom}_{\mathcal H}(e_{\mathcal H},e_{\mathcal H}\times I)
\]
% margin note written sideways
\margin{\unsure{Ça a tendance à} être \unsure{archi-faux} si $\mathcal H$ est additive. Par contre, vrai si $\mathcal H$ est une catégorie à sommes distributives \unsure{p.r.|p/r} aux produits (voire universelles) dont l'objet final est \uline{connexe} \ldots\ $\rightarrow$}
soit \uline{bijective},\inserted{$\ast$ i.e. que $I\to I_W$ soit \uline{pl. fid.}} On définit alors, pour \unsure{tt} $I$, une application canonique
\diagram{\begin{tikzcd}[ampersand replacement=\&,column sep=small]
\mathrm{Hom}(X,I_{\mathcal H}) \arrow[rr] \arrow[dr,dashed,"{\pi_0^{X,I_{\mathcal H}}}"'] \& \& \mathrm{Hom}(\pi_0(X),I) \arrow[dl,dashed,"\sim"]\\
\& \mathrm{Hom}(\pi_0(X),\pi_0(I_{\mathcal H})) \& 
\end{tikzcd}.}

\deleted{\unsure{L}}\unsure{La} ``bi-adjonction'' des foncteurs $\pi_0$ et $\unsure{\mathcal H|I}\mapsto I_{\mathcal H}$ revient à dire que cette flèche est \deleted{\unsure{p}}\unsure{bijective},
